% Generated draft: requires mathematical and source review.
\documentclass[11pt]{article}
\usepackage[margin=25mm]{geometry}
\usepackage{fontspec}
\setmainfont{Latin Modern Roman}
\usepackage{amsmath,amssymb,mathtools,unicode-math}
\setmathfont{Latin Modern Math}
\usepackage{enumitem,needspace,xurl,hyperref}
\setlength{\parindent}{0pt}
\setlength{\parskip}{5pt}
\setlength{\emergencystretch}{3em}
\newcommand{\N}{\mathbb{N}}
\newcommand{\Z}{\mathbb{Z}}
\newcommand{\Q}{\mathbb{Q}}
\newcommand{\R}{\mathbb{R}}
\newcommand{\C}{\mathbb{C}}
\newcommand{\F}{\mathbb{F}}
\newcommand{\A}{\mathbb{A}}
\newcommand{\PP}{\mathbb{P}}
\newcommand{\E}{\mathbb{E}}
\newcommand{\eps}{\varepsilon}
\newcommand{\dd}{\,\mathrm{d}}
\begin{document}
\section*{UnsolvedMath ID: 30000648}
\textbf{Source problem code:} catalogue entry
\subsection{Definitions and mathematical statement}
For each positive integer $n$, work in $\C[X_1,\ldots,X_n]$ and define the constant-coefficient differential operator
\[
 \Delta=\sum_{i=1}^n\frac{\partial^2}{\partial X_i^2}.
\]
Let $P$ be homogeneous of degree four. The notation $\Delta^m$ means composition of $\Delta$ with itself $m$ times, while $P^m$ means ordinary multiplication of the polynomial. Every vanishing equality below is an identity in the polynomial ring.

The Vanishing Conjecture in the source asks whether
\[
 \left[\forall m\in\N_{>0}:\quad\Delta^m(P^m)=0\right]
 \quad\Longrightarrow\quad
 \left[\exists M\in\N_{>0}\ \forall m\ge M:\quad
                      \Delta^m(P^{m+1})=0\right]
\]
for every such $P$ and every number of variables. The threshold $M$ may depend on $P$ and $n$; no uniform numerical bound is requested. The conclusion requires vanishing at every sufficiently large integer, not only along a subsequence. The hypothesis is the full collection of identities for all positive $m$, not merely harmonicity $\Delta P=0$. The same exponent $m$ applies to the differential operator in both clauses; only the polynomial power changes from $m$ to $m+1$.
\subsection{Short English statement}
For a homogeneous quartic over $\C$, does vanishing of every iterated-Laplacian expression $\Delta^m(P^m)$ force $\Delta^m(P^{m+1})$ to vanish for all sufficiently large $m$? The eventual threshold may depend on the polynomial.
\subsection{Sources}
\begin{itemize}
\item Local UnsolvedMath registry, permanent ID 30000648, code catalogue entry, ``Vanishing Conjecture for Powers of Homogeneous Polynomials.'' Proposed year 2007; proposer unspecified. The original statement labels this the Vanishing Conjecture.
\item \emph{Affine Algebraic Geometry} (2007), Oberwolfach Reports, pp.\ 5--82, DOI 10.4171/owr/2007/01; \url{https://doi.org/10.4171/owr/2007/01}. This is the sole supplied literature reference.
\item Statement status corrected\_verified; the supplied note says the source label was removed and operator, ring, degree, powers, and quantifiers were enclosed in balanced delimiters.
\item Literature triage checked 2026-08-21: open. It generically describes the item as unresolved and says no authoritative complete resolution was found by the supplied targeted check; it provides no specific partial cases or theorem claims. Created 2026-08-25; updated 2026-09-16. Current status and source scope require independent verification.
\item Registry curation by Ulam AI / UnsolvedMath Contributors, CC BY 4.0; the underlying report retains its supplied provenance. Unreviewed statement-only draft.
\end{itemize}
\end{document}
